\documentclass[11pt]{article}

\usepackage[margin=1in]{geometry}
\usepackage[T1]{fontenc}
\usepackage[utf8]{inputenc}
\usepackage[expansion=false]{microtype}
\usepackage{booktabs}
\usepackage{enumitem}
\usepackage{xcolor}
\usepackage{hyperref}

\newcommand{\todo}[1]{\textcolor{red}{[TODO: #1]}}
\newcommand{\project}{\textsc{Adventure-2600}}

\title{\project: A Protocol for Camera-Mediated Agent Evaluation on an Original Atari 2600}
\author{Anonymous authors\\
\texttt{\todo{replace with author names, affiliations, and contact}}}
\date{\todo{replace with submission date}}

\begin{document}
\maketitle

\begin{abstract}
Atari benchmarks have played an influential role in reinforcement-learning
research, but most reported evaluation is performed through software
environments with access to convenient resets, frame buffers, and score
signals.  We present a preregistered experimental protocol for evaluating an
agent on an unmodified original Atari 2600.  The agent is trained in an
emulator and then plays \emph{Adventure} by observing a physical game display
through a fixed USB camera and by issuing only standard joystick switch
closures through an external controller bridge.  The system records
synchronized video, requested actions, applied relay states, and timing so
that behaviour can be inspected and replayed.  This manuscript defines the
hardware boundary, transfer procedure, baselines, metrics, and reporting
requirements; it reports no empirical performance results.  It neither seeks
to reproduce Agent57's Atari57 result nor to claim general or human-level
intelligence.  The objective is to make a future real-hardware result
falsifiable and reproducible.
\end{abstract}

\section{Introduction}

The Arcade Learning Environment (ALE) made Atari 2600 games a common platform
for studying general-purpose reinforcement-learning algorithms
\cite{bellemare2013arcade}.  Deep Q-Networks and their descendants established
strong Atari results \cite{mnih2015human,kapturowski2019recurrent}, and
Agent57 emphasized breadth across the 57-game suite \cite{badia2020agent57}.
These software benchmarks are valuable, but their interfaces are not the same
as an agent operating a physical console: a deployed system must cope with
video capture, display/camera distortion, latency, external actuation, and
less convenient episode control.

This paper scopes a narrow real-hardware experiment.  We ask whether an agent
trained in an Atari emulator can transfer to \emph{Adventure} on an original
Atari 2600 while being evaluated from camera pixels and standard joystick
inputs only.  The work is deliberately not an attempt to establish broad
generality.  Instead, it makes the systems boundary and evidence requirements
explicit before results are collected.

\subsection{Agent57 as motivation, not a target to reproduce}

Agent57 is the most directly relevant reference point: it trained a family of
policies ranging from exploratory to exploitative and adaptively prioritized
them, reporting the first deep-RL result above the standard human benchmark on
all 57 Atari games \cite{badia2020agent57}.  Its emphasis on weak-task
performance and the exploration--exploitation trade-off is a useful design
lesson for our proposed memory and exploration components.

Our study is complementary rather than a like-for-like comparison.  Agent57's
claim concerns a software Atari benchmark; its evaluation interface supplies
an emulator's observations, actions, resets, and learning signals.  It does
not test the additional physical-system questions introduced here: camera and
display variation, externally applied controller contacts, end-to-end latency,
and an audit trail that lets an observer see what was applied to the console.
Those are not defects in Agent57---they were outside its stated benchmark
scope---but they delimit what its reported score can establish about deployed
physical play.

We therefore challenge three narrower assumptions rather than Agent57's
Atari57 result: (i) that high simulator performance is sufficient evidence of
real-console transfer; (ii) that average or normalized score alone is an
adequate report when hardware timing and failures matter; and (iii) that the
large-scale, adaptive-policy approach is the right efficiency baseline for a
single accessible physical setup.  We will compare our method only with
declared same-budget baselines, report all failures and compute, and never
interpret a result on one game as evidence of general Atari competence.

\section{Research questions and hypotheses}

\paragraph{RQ1: transfer.} Can an emulator-trained agent achieve a predefined
\emph{Adventure} objective on original hardware using only camera observations
and joystick actions?  The primary hypothesis is that latency- and
appearance-randomized training improves zero-shot transfer over a baseline
trained with a fixed, idealized emulator interface.

\paragraph{RQ2: efficiency.} At a fixed interaction-frame budget, does the
proposed memory/exploration design improve the selected task metric relative to
a recurrent pixel-policy baseline?  The study will report interaction frames,
wall-clock training time, and compute/energy estimates rather than only final
score.

\paragraph{RQ3: observability.} Can every hardware episode be reconstructed
from synchronized camera video, requested actions, applied relay actions, and
configuration metadata?  The hypothesis is operational: a third party can
verify alignment and reproduce the reported outcome from the released log.

\section{Experimental system}

\subsection{Environment and task}

The initial task is Atari 2600 \emph{Adventure}, selected because its
standard joystick/fire interface combines navigation, object interaction,
delayed outcomes, and room-to-room memory.  Before collecting results, we
will declare the ROM hash, game variation, reset/start procedure, console
region, and a finite set of task endpoints.  Candidate endpoints include
reliable movement, carrying an object across a room transition, returning to a
previously observed room, and task completion.  The final endpoint definitions
and scoring rules must be fixed before the evaluation runs are inspected.

\subsection{Emulator training}

The emulator API exposes \texttt{reset(rom, seed)} and
\texttt{step(action, repeat\_frames)}.  During training, reward and terminal
signals may be used for learning; the physical-console evaluation policy does
not receive score, RAM, room identifiers, object coordinates, or emulator
state.  The baseline is a recurrent pixel policy with prioritized replay.  The
proposed agent augments it with a compact latent world model, uncertainty-aware
exploration, and conditional planning.  Model capacity, optimizer, replay
policy, seeds, action repeat, and total interaction budget will be disclosed.

To prepare for transfer, training randomizes camera-relevant appearance and
timing: random crops and perspective perturbations, colour and brightness
changes, skipped frames, action holds, and measured camera/display/control delay.  Randomization
ranges will be chosen from calibration data, not tuned against final hardware
scores.

\subsection{Original-console interface}

Figure~\ref{fig:system} is a placeholder for the final system diagram.  An
original Atari 2600 drives its normal display.  A fixed USB camera mounted on
a two-axis Robotis head observes that display; pan and tilt are used only for
setup and then fixed.  An existing Linux box hosts camera capture, policy
inference, preview, and logging.  It transmits bounded commands over USB
serial to a classic Arduino Nano.  The Nano validates legal joystick masks
and, through a ULN2803A driver and five normally-open relays, closes the
standard up, down, left, right, and fire lines to controller common.  The
relay contacts are the sole electrical connection from the agent apparatus to
the controller port.

\begin{figure}[t]
\centering
\fbox{\parbox{0.9\linewidth}{\centering
\todo{System figure: Atari $\rightarrow$ display $\rightarrow$ fixed camera
$\rightarrow$ Linux box $\rightarrow$ Nano $\rightarrow$ relay contacts
$\rightarrow$ Atari controller port.  Label the synchronized log.}}}
\caption{Planned physical-console evaluation boundary.}
\label{fig:system}
\end{figure}

The original joystick remains unmodified and is used separately for human
reference runs.  It is never connected in parallel with the relay bridge.
Console power, reset, difficulty, and game selection remain manual in this
experiment.  A direct composite-to-USB capture alternative may be evaluated
later but is not part of the primary configuration.

\subsection{Safety and failure behaviour}

All relay contacts must start open.  The Nano opens every contact on boot,
malformed commands, serial loss, or watchdog timeout; it rejects opposing
directions and bounds a hold duration unless the host renews it.  Hardware is
commissioned with continuity tests before console connection.  No controller
port pin powers relay coils or receives a driven host voltage.

\section{Evaluation protocol}

\subsection{Conditions}

We will compare the following conditions using the same declared game
configuration:
\begin{enumerate}[leftmargin=*]
  \item recurrent emulator baseline, evaluated in the emulator;
  \item proposed agent, evaluated in the emulator;
  \item frozen proposed agent, evaluated on original hardware;
  \item where justified by a preregistered budget, proposed agent with limited,
        documented hardware fine-tuning; and
  \item human reference runs using the original joystick.
\end{enumerate}
Human runs are descriptive references, not training demonstrations and not a
claim of statistical equivalence without a preregistered protocol.

\subsection{Metrics}

The primary metric will be the preregistered task endpoint success rate across
held-out seeds/starts, with binomial confidence intervals.  Secondary metrics
are time to endpoint, score where lawfully observable and separately recorded,
survival/progress measures, emulator interaction frames, wall-clock time,
compute/energy, and median/90th-percentile/worst-case end-to-end latency.
We will report all attempted hardware episodes, including aborts, watchdog
trips, dropped-frame runs, and failures to start.

\subsection{Observability and audit trail}

Each episode produces camera video or timestamped frames, requested action,
applied relay mask, monotonic timestamps, configuration hash, and outcome.
A playback tool overlays actions and timing on recorded gameplay.  We will
test synchronization by issuing repeated visible actions and measuring the
time from host command through relay application to the first responding camera
frame.  The released artifact must permit an independent reviewer to inspect
behaviour without access to hidden console or emulator state.

\section{Analysis plan}

We will define the sample size, seeds/starts, exclusion rules, randomization
ranges, and confidence-interval method before collecting final hardware
results.  Hyperparameter selection occurs on disjoint development seeds and
may not use held-out hardware outcomes.  We will publish learning curves and
per-run results rather than only maxima, and distinguish zero-shot transfer
from any hardware-adapted condition.  Claims will be limited to the declared
task, hardware configuration, and compute budget.

\section{Limitations and responsible reporting}

This single-game study cannot establish general intelligence, broad Atari
competence, or human-level performance.  A camera-plus-display path introduces
latency and visual artefacts that are absent from emulator pixels; conversely,
it is a useful test of external perception and action.  Original-console
variation, display properties, ROM rights, and human-reference methodology
must be reported.  We will release only material we are entitled to share and
will avoid distributing copyrighted ROMs.

\section{Reproducibility checklist}

\begin{itemize}[leftmargin=*]
  \item ROM hash and lawful acquisition procedure (without redistributing ROMs).
  \item Console, display, camera, Linux host, Nano, relay and firmware revisions.
  \item Camera crop, exposure, gain, white balance, frame rate, and Robotis
        head positions.
  \item Relay wiring diagram, continuity tests, watchdog tests, and command
        protocol.
  \item Training code, environment version, seeds, budgets, checkpoints, and
        dependency lockfile.
  \item Raw/re-encoded gameplay logs, action/relay traces, latency calibration,
        analysis code, and complete result tables.
\end{itemize}

\section{Conclusion}

This draft specifies a modest and auditable experiment: an agent learns in an
emulator, plays \emph{Adventure} on an original Atari 2600 through normal
human-style controls, and leaves an inspectable record of what it saw and did.
Empirical conclusions are intentionally deferred until the specified evidence
has been collected.

\bibliographystyle{plain}
\bibliography{references}

\end{document}
